\chapter{Discussion}
\label{cha:discussion}

\section{Variance Analysis}

\subsection{Baseline MC Estimator}

The MC estimator draws $N$ samples ${\xm}^{(i)} \sim p(\xm \mid \xo)$ and returns
\[
  \widehat{p}_{\mathrm{MC}}(N) = \frac{1}{N}\sum_{i=1}^{N} \sigma\!\left(w^T x^{(i)}\right).
\]
Since $\sigma(\cdot) \in (0,1)$, the per-sample variance $\sigma^2_{\mathrm{MC}}
:= \mathrm{Var}_{p(\xm\mid\xo)}\!\left[\sigma(w^T x)\right]$ is bounded by $1/4$,
giving
\[
  \mathrm{Var}\!\left[\widehat{p}_{\mathrm{MC}}(N)\right]
  = \frac{\sigma^2_{\mathrm{MC}}}{N}
  \;\leq\; \frac{1}{4N}.
\]

\subsection{PCPG-MC Estimator}

The PCPG-MC estimator with $M$ outer (PG) and $L$ inner (Gaussian) samples returns
\[
  \widehat{p}_{\mathrm{PCPG}}(M,L)
  = \frac{1}{2ML}\sum_{j=1}^{M}\sum_{i=1}^{L}
    \mathrm{Re}\!\left[h(\gamma_j, u_{ij})\right],
\]
where
$h(\gamma, u) := \exp\!\left(\tfrac{1}{8\gamma} + iu\!\left(\tfrac{1}{2\gamma}-c\right)\right)
\cdot \phi_{\xm\mid\xo}(-uw_m)$,
$c = w_o^T\xo$, and $u_{ij} \sim \mathcal{N}(0,\gamma_j)$.

\paragraph{Output range.}
A single sample ($M{=}L{=}1$) returns $\tfrac{1}{2}\mathrm{Re}[h(\gamma,u)]$.
Since $|\phi| \leq 1$ and $|h(\gamma,u)| = \exp\!\left(\tfrac{1}{8\gamma}\right)$,
the value lies in $\left[-\tfrac{1}{2}e^{1/(8\gamma)},\, +\tfrac{1}{2}e^{1/(8\gamma)}\right]$.
Because $\gamma \sim \PG(1,0)$ can be arbitrarily small, $e^{1/(8\gamma)}$ is
unbounded, so individual PCPG samples are not confined to $[0,1]$ and can be
negative or exceed 1.
Empirically, with $M{=}1000$, $L{=}1$ the standard deviation exceeds $0.76$,
and with $M{=}10{,}000$, $L{=}1$ it exceeds $1.38$.
By contrast, every individual MC sample lies in $(0,1)$.

\paragraph{Variance decomposition.}
Define the inner conditional expectation
$g(\gamma) := \E_{u \sim \mathcal{N}(0,\gamma)}[\mathrm{Re}\,h(\gamma, u)]$
and two variance components:
\[
  \sigma^2_B := \mathrm{Var}_{\gamma \sim \PG(1,0)}\!\left[g(\gamma)\right]
  \quad\text{(outer)}, \qquad
  \sigma^2_A := \E_{\gamma}\!\left[\mathrm{Var}_{u|\gamma}\!\left[\mathrm{Re}\,h(\gamma,u)\right]\right]
  \quad\text{(inner)}.
\]
By the law of total variance:
\[
  \mathrm{Var}\!\left[\widehat{p}_{\mathrm{PCPG}}(M,L)\right]
  = \frac{1}{4}\!\left(\frac{\sigma^2_B}{M} + \frac{\sigma^2_A}{ML}\right).
\]
These two components behave very differently. $\sigma^2_B$ is \emph{finite}:
evaluating the Gaussian integral in closed form shows that $g(\gamma)$ remains
$O(1)$ as $\gamma \to 0$ because the $e^{1/(8\gamma)}$ prefactor in $h$ is
exactly cancelled by the saddle-point of the Gaussian integral. By contrast,
$\sigma^2_A$ is \emph{infinite}: for fixed $\gamma$,
$\mathrm{Var}_{u|\gamma}[\mathrm{Re}\,h] \approx \tfrac{1}{2}e^{1/(4\gamma)}$,
and the PG density decays only as $\gamma^{-3/4}e^{-1/(8\gamma)}$ near zero,
so $\E_\gamma[e^{1/(4\gamma)}]$ diverges.

\paragraph{Optimal budget split.}
Because $\sigma^2_A = \infty$, the formula above breaks down for $L = 1$:
individual samples $\mathrm{Re}\,h(\gamma, u)$ have unbounded variance due to
the unmitigated $e^{1/(8\gamma)}$ factor. The inner average $(1/L)\sum_i
\mathrm{Re}\,h(\gamma, u_i)$ approaches the finite $g(\gamma)$ as $L$
grows, stabilising the estimator. Therefore, for a fixed budget $N = ML$,
\emph{more Gaussian inner samples $L$ is better, not more PG samples $M$}.
This is confirmed empirically: at $N{=}10{,}000$ the standard deviation
falls from $1.38$ ($M{=}10{,}000$, $L{=}1$) to $0.04$ ($M{=}10$,
$L{=}1{,}000$). The right strategy is to keep $M$ moderate and let $L$ be
large. PCPG-GH takes this to its logical conclusion by making $L \to \infty$
(exact quadrature) for each PG sample.

\subsection{PCPG-GH Estimator}

Replacing the inner MC with $N_H$-point Gauss--Hermite quadrature
eliminates $\sigma^2_A$ entirely. For a fixed number of PG samples $M$:
\[
  \mathrm{Var}\!\left[\widehat{p}_{\mathrm{GH}}(M)\right]
  = \frac{\sigma^2_B}{4M}.
\]
Compared to PCPG-MC with the optimal split ($L{=}1$, same total budget
$N = M \cdot N_H$ evaluations), PCPG-GH uses $M = N/N_H$ PG samples but
removes $\sigma^2_A$:
\[
  \mathrm{Var}\!\left[\widehat{p}_{\mathrm{GH}}\right]
  = \frac{N_H\,\sigma^2_B}{4N}
  \quad\text{vs}\quad
  \mathrm{Var}\!\left[\widehat{p}_{\mathrm{PCPG}}\right]
  = \frac{\sigma^2_B + \sigma^2_A}{4N}.
\]
PCPG-GH is better when $\sigma^2_A > (N_H - 1)\,\sigma^2_B$, i.e.\ when the
inner variance is large relative to the outer. In practice, because the
integrand oscillates heavily for small $\gamma$ (making $\sigma^2_A$ large),
PCPG-GH tends to substantially reduce overall variance compared to PCPG-MC.
The outer variance $\sigma^2_B/M$ remains and still exceeds the MC bound
$\sigma^2_{\mathrm{MC}}/N$.

\subsection{PCPG-Quad Estimator}

Replacing both integrals with quadrature (Golub--Welsch for the outer PG
integral, Gauss--Hermite for the inner) makes the estimator fully deterministic:
\[
  \mathrm{Var}\!\left[\widehat{p}_{\mathrm{Quad}}\right] = 0.
\]
The estimator always returns the same value for the same inputs, and its
output is bounded since the PG quadrature nodes $\gamma_j$ are fixed
positive numbers --- the factor $e^{1/(8\gamma_j)}$ is a known finite
constant for each node.
The error is purely from quadrature approximation (bias), not variance.

\section{Bias Analysis}

\subsection{MC Estimator}
MC is unbiased for all $N \geq 1$: sampling from $p(\xm \mid \xo)$ and
averaging $\sigma(w^T x)$ is a direct Monte Carlo estimate of the target
expectation, so $\E[\widehat{p}_{\mathrm{MC}}(N)] = p(y \mid \xo)$ exactly.

\subsection{PCPG-MC Estimator}
PCPG-MC is also unbiased. Both integrals are approximated by ordinary MC,
so by linearity of expectation the estimator equals the true double integral,
which is $p(y \mid \xo)$ by the PG-sigmoid and Hubbard--Stratonovich
identities.

\subsection{PCPG-GH Estimator}
The outer PG integral is still MC (unbiased). The inner integral is replaced
by Gauss--Hermite quadrature, which introduces a small approximation error.
For $N_H$ quadrature nodes, GH integrates polynomials of degree up to $2N_H - 1$
exactly, and for smooth (analytic) integrands the error decays
super-algebraically in $N_H$. In our setting the integrand is analytic in $u$,
so the bias is negligible already for small $N_H$ (in practice $N_H \geq 20$
is sufficient).

\subsection{PCPG-Quad Estimator}
Both integrals introduce a quadrature error. The Gauss--Hermite inner error
is the same as for PCPG-GH. The Golub--Welsch PG quadrature with $M$ nodes
integrates polynomials in $\gamma$ exactly up to degree $2M-1$; for smooth
functions of $\gamma$ the error again decays rapidly with $M$. The maximum
reliable order is $M \leq 19$ (limited by numerical precision in the moment
computation). At $M = 19$ and $N_H \geq 20$, both errors are numerically
negligible and the output closely matches the unbiased MC ground truth.

\section{Comparison}

Table~\ref{tab:estimator-comparison} summarises the properties of all four
estimators.

\begin{table}[H]
\centering
\small
\begin{tabular}{lllll}
\hline
\textbf{Estimator} & \textbf{Bias} & \textbf{Variance} & \textbf{Output range} & \textbf{Deterministic} \\
\hline
MC           & none  & $\sigma^2_{\mathrm{MC}}/N \leq 1/(4N)$ & $(0,1)$    & no  \\
PCPG-MC      & none  & $(\sigma^2_B/M + \sigma^2_A/ML)/4$     & unbounded  & no  \\
PCPG-GH      & small & $\sigma^2_B/(4M)$                      & unbounded  & no  \\
PCPG-Quad    & small & $0$                                    & fixed      & yes \\
\hline
\end{tabular}
\caption{Summary of estimator properties. $\sigma^2_{\mathrm{MC}} \leq 1/4$ is bounded; $\sigma^2_B$ and $\sigma^2_A$ are not.}
\label{tab:estimator-comparison}
\end{table}

\paragraph{Variance.}
MC has the best variance guarantee: $\sigma^2_{\mathrm{MC}} \leq 1/4$ provides
a tight universal bound, independent of the model or weights. The PCPG
estimators have no such bound --- both $\sigma^2_B$ and $\sigma^2_A$ depend
on the problem and can be large. In the experiments, PCPG-MC with an equal
split yields $\sim$26$\times$ higher standard deviation than MC at the same
budget. PCPG-GH narrows but does not close this gap.
PCPG-Quad has zero variance by construction.

\paragraph{Bias.}
MC and PCPG-MC are exactly unbiased. PCPG-GH and PCPG-Quad introduce small
quadrature errors that are negligible at the parameter choices used
($N_H \geq 20$, $M_{\mathrm{PG}} \leq 19$).

\paragraph{When to use which.}
MC is the simplest and most reliable choice when the density $p(\xm \mid \xo)$
is easy to sample from and low variance is the priority.
PCPG-MC becomes preferable when sampling is difficult but evaluating the
characteristic function of $p(\xm \mid \xo)$ is cheap (as with a GMM), since
PCPG never samples the missing features directly.
PCPG-GH is a direct improvement over PCPG-MC: it eliminates inner variance
at no extra cost for the same outer budget, and should be preferred whenever
$\sigma^2_A$ is non-negligible.
PCPG-Quad is the right choice when a fully deterministic, reproducible result
is required --- for example when differentiating through the estimator or when
multiple evaluations at the same point must agree exactly.


\section{Effect of Dimensionality}

The number of missing features $d_m$ affects the two families of estimators in
opposite ways.

\paragraph{MC.}
The per-sample variance $\sigma^2_{\mathrm{MC}} = \mathrm{Var}[\sigma(w^T x)]$
grows with $d_m$: more missing features make $w^T x$ more random, spreading the
distribution of $\sigma(w^T x)$ and pushing $\sigma^2_{\mathrm{MC}}$ toward its
upper bound of $1/4$. The variance of the MC estimator therefore increases with
dimensionality.

\paragraph{PCPG.}
The situation is reversed. For a GMM with diagonal covariance the characteristic
function at frequency $s = -u w_m$ contains the damping factor
\[
  \left|\varphi(-u w_m \mid \xo)\right|
  \;\leq\; \exp\!\left(-\tfrac{1}{2}u^2 \sum_{j \in \mathrm{mis}} w_j^2 \sigma_{kj}^2\right),
\]
where the exponent grows with $d_m$. This means the CF decays faster in $u$ as
more features are missing, concentrating the integrand near $u = 0$ where the
$e^{1/(8\gamma)}$ blowup from small $\gamma$ is suppressed. As a result,
$\sigma^2_A$ (the inner variance) decreases with $d_m$, and $\sigma^2_B$ is also
damped. The PCPG estimators become \emph{more} accurate, not less, as
dimensionality grows.

The same damping also reduces the Gauss--Hermite quadrature error for PCPG-GH:
a narrower integrand is easier to approximate with a fixed number of nodes, so
the bias of PCPG-GH decreases with $d_m$ as well.

\paragraph{Summary.}
MC variance increases toward $1/(4N)$ as $d_m$ grows, while PCPG variance
decreases due to CF concentration. The relative advantage of PCPG over MC
therefore improves with the number of missing features, which is precisely the
high-missing-data regime where accurate marginalisation matters most.

\section{Related Work}

Here, we just focus on works that also try to calculate 

\[
  \mathbb{E}_{P(\xm \mid \xo)} \left[\sigma(w^T x)\right]
\]

but use a different approach.

\subsection{Naive Conformant Learning (NaCL)}

\citet{khosravi2019expect} first point out that this expectation is hard
to compute in general: even with a naive Bayes feature distribution, it
is NP-hard. They also note that the popular shortcut of mean imputation
only matches the expectation when the classifier is linear \emph{and} the
features are independent --- for the sigmoid it is always biased.

Their fix is to pick a feature distribution on which the expectation
\emph{is} cheap, namely a naive Bayes model over binary features. The
trick is that naive Bayes and logistic regression are dual to each other:
every naive Bayes induces a logistic regression, and many naive Bayes
models induce the same one. They call a naive Bayes \textit{conformant}
with the given LR if both give the same class probabilities on fully
observed inputs, and learn such a distribution via geometric programming
(``naive conformant learning''). Once this naive Bayes is fitted, the
expected prediction is just a marginal inference query and runs in linear
time.

\medskip
\noindent\textbf{Difference to us:} they keep the feature distribution
simple (naive Bayes on binary inputs) so the expectation is exact and
cheap; we keep the feature distribution rich (a GMM or PC on continuous
inputs) and instead use a stochastic trick to handle the sigmoid. They
also have to refit the feature distribution whenever the classifier
changes, while ours is trained once and reused.

\subsection{Joint Modelling via SAEM}

\citet{jiang2020logreg} assume a joint Gaussian model on the covariates,
$x \sim \mathcal{N}(\mu, \Sigma)$, together with logistic regression on
top. Their main contribution is on the training side: they fit
$(\mu, \Sigma, \beta)$ jointly from training data that may itself contain
missing values, using a stochastic-approximation EM algorithm with
Metropolis--Hastings updates.

For prediction at test time, they simply do plain Monte Carlo: sample
$x^{m,(s)} \sim p(\xm \mid \xo)$ from the fitted Gaussian and average,
\[
  p(y \mid \xo) \approx \frac{1}{S} \sum_{s=1}^{S} p\!\left(y \mid \xo, x^{m,(s)}\right)
\]
This is exactly the MC baseline we compare against, just with a single
Gaussian instead of a mixture as the density.

\medskip
\noindent\textbf{Difference to us:} they sample the missing features and
average the sigmoid; we never sample the missing features, but instead
rewrite the sigmoid so that only the characteristic function of
$p(\xm \mid \xo)$ is needed. Their feature distribution is also limited
to a single Gaussian, while ours can capture multimodal shapes.
